% SEGMENT OLP-0603-S01
% Part: methods
% Chapter: proofs
% Section: introduction

\documentclass[../../../include/open-logic-section]{subfiles}

\begin{document}

\olfileid{mth}{prf}{int}

\olsection{அறிமுகம்}

அறிமுகத் தருக்கத்தில் உங்களுக்கு ஏற்பட்ட அனுபவத்தால்,
!!a{derivation} முறைமை உங்களுக்குப் பழக்கமானதாக இருக்கலாம்.
அது இயற்கை வருவித்தல் முறைமையாகவோ ஃபிட்ச் பாணி
!!{derivation} முறைமையாகவோ, ஒருவேளை நிறுவல்-மர
முறைமையாகவோ இருக்கலாம். இம்முறைமைகளில் நீங்கள் நிறுவல்களைச்
செய்தது நினைவிருக்கலாம்: !!a{formula}வை நிறுவுவது அல்லது
தரப்பட்ட வாதம் செல்லுபடியாகும் என்று காட்டுவது. இதற்காக,
விரும்பிய இறுதி முடிவைப் பெறும்வரை முறைமையின் விதிகளைப்
பயன்படுத்தினீர்கள். தருக்கத்தைப் \emph{பற்றி} ஆராயும்போதும்
நாம் கருத்துகளை நிறுவுகிறோம்; ஆனால் பெரும்பாலும்
!!a{derivation} முறைமையைப் பயன்படுத்துவதில்லை. நாம் கருதும்
பெரும்பாலான நிறுவல்கள் முதல்தர தருக்க மொழியிலேயே முழுவதுமாக
இருப்பதில்லை; ஆங்கிலத்தில் எழுதப்படுகின்றன, இடையிடையே குறியீட்டு
மொழியும் வரலாம். இப்படிப்பட்ட நிறுவலை உருவாக்கத் தொடங்கும்போது,
!!a{derivation} முறைமை இல்லாமல் எவ்வாறு நிறுவுவது, எங்கிருந்து
தொடங்குவது, தமது நிறுவல் சரியானதா என்று எப்படித் தெரிந்துகொள்வது
எனத் திகைக்கலாம்.

% SEGMENT OLP-0603-S02
ஒரு நிறுவலை முயல்வதற்கு முன், அது என்ன என்றும் அதை எப்படி
உருவாக்குவது என்றும் அறிந்துகொள்ள வேண்டும். பெயர் உணர்த்துவதுபோல்,
ஒரு \emph{நிறுவல்} ஏதோ ஒன்று மெய் என்பதைக் காட்டுகிறது. இதை
ஓர் உரையாடலாகக் கருதலாம்: இரண்டு தவிர்ந்த ஒவ்வொரு பகா எண்ணும்
ஒற்றை எண்ணா என்று ஒருவர் கேட்கிறார் என்க. ``ஆம்'' என்று மட்டும்
பதிலளிப்பது போதாது; அது \emph{ஏன்} என்று அவர் அறிய விரும்பலாம்.
அப்போது நீங்கள் அவருக்கு ஒரு நிறுவலைத் தருவீர்கள்.

% SEGMENT OLP-0603-S03
அன்றாட உரையாடலில் ஒரு விடையைக் குறிப்பாகச் சுட்டுவதோ
முழுமையற்ற விடை தருவதோ போதுமானதாக இருக்கலாம். ஆனால்
தருக்கத்திலும் கணிதத்திலும் கண்டிப்பான நிறுவல் வேண்டும்:
\emph{எந்த} ஐயத்துக்கும் இடமின்றி ஒன்று மெய் என்று காட்ட
வேண்டும். ஆகவே நிறுவலின் ஒவ்வொரு படிக்கும் நியாயம் இருக்க
வேண்டும்; அந்த நியாயமும் பொருத்தமானதாக இருக்க வேண்டும்.
எடுத்துக்காட்டாக, நீங்கள் பயன்படுத்தும் எடுகோள் நிறுவப்படும்
தேற்றத்தின் கூற்றிலேயே உண்மையாக எடுக்கப்பட்டிருக்க வேண்டும்;
வரையறைகளைச் சரியாகப் பயன்படுத்த வேண்டும்; நியாயங்களுக்காகச்
சுட்டப்படும் உய்த்தறிதல்கள் சரியானவையாக இருக்க வேண்டும்.

% SEGMENT OLP-0603-S04
வழக்கமாக நாம் ஒரு கூற்றை நிறுவுகிறோம். நிறுவப்படும் கூற்றுகளை
முன்மொழிவுகள், தேற்றங்கள், துணைத்தேற்றங்கள், தொடர்விளைவுகள்
என்று பல பெயர்களால் அழைக்கிறோம். முன்மொழிவு என்பது
நிறுவி பதிவுசெய்யும் அளவு முக்கியமான, ஆனால் மிக ஆழமானதாகவோ
அடிக்கடி பயன்படுவதாகவோ இருக்க வேண்டியதில்லை என்ற கூற்று.
தேற்றம் என்பது குறிப்பிடத்தக்க, முக்கியமான முன்மொழிவு. அதன்
நிறுவல் பல படிகளாகப் பிரியலாம். சில நேரம் அதை முதலில்
நிறுவியவரின் பெயரால், எடுத்துக்காட்டாக காண்டோரின் தேற்றம்
அல்லது L\"owenheim--Skolem தேற்றம் என, அழைப்பார்கள்;
சில நேரம் அது பற்றிய முடிவின் பெயரால், எடுத்துக்காட்டாக
நிறைவுத் தேற்றம் என, அழைப்பார்கள். துணைத்தேற்றம் என்பது
மேலும் முக்கியமான முடிவின் நிறுவலில் பயன்படுத்தப்படும்
முன்மொழிவோ தேற்றமோ ஆகும். குழப்பமூட்டும் விதமாகச் சில
துணைத்தேற்றங்கள் தாமே முக்கியமான முடிவுகளாக இருக்கும்;
அவற்றை அறிமுகப்படுத்தியவரின் பெயரையும் பெறும், எடுத்துக்காட்டாக
சோர்னின் துணைத்தேற்றம். மற்றொரு முடிவிலிருந்து எளிதில்
பின்தொடரும் முடிவே தொடர்விளைவு.

% SEGMENT OLP-0603-S05
நிறுவ வேண்டிய கூற்றில், எப்படிப்பட்ட பொருள்களைப் பற்றிப்
பேசுகிறோம் என்பதைத் தெளிவாக்கும் எடுகோள்கள் அடிக்கடி இருக்கும்.
``$!B \lif !C$ வடிவிலுள்ள !!a{formula}வாக $!A$ இருக்கட்டும்''
என்றோ, ``$\Gamma \Proves !A$ என்க'' என்றோ அது தொடங்கலாம்.
இவை முன்மொழிவின், தேற்றத்தின் அல்லது துணைத்தேற்றத்தின்
\emph{எடுகோள்கள்}; நிறுவலில் இவற்றை மெய்யாக எடுத்துக்கொள்ளலாம்.
இவை நாம் நிறுவுவதன் பரப்பைக் கட்டுப்படுத்துவதோடு, பேசப்படும்
பொருள்களுக்குச் சில பெயர்களையும் அறிமுகப்படுத்துகின்றன.
எடுத்துக்காட்டாக, ``$!B \lif !C$ வடிவிலுள்ள !!a{formula}வாக
$!A$ இருக்கட்டும்'' என்று உங்கள் முன்மொழிவு தொடங்கினால்,
ஒரு குறிப்பிட்ட வகை வாய்பாடுகளைப் பற்றி மட்டுமே, அதாவது
நிபந்தனை வாய்பாடுகளைப் பற்றி மட்டுமே, நிறுவுகிறீர்கள்.
உங்கள் நிறுவலில் பேசப்படும் $!B \lif !C$ ஒரு தன்னிச்சையான
நிபந்தனை வாய்பாடு என்பது அப்போது புரிந்துகொள்ளப்படுகிறது.

\end{document}
